\documentclass[12pt]{article}

\usepackage{amsmath, geometry, amssymb, bm, graphicx}
\geometry{margin=1in}
\newcommand{\dd}{\mathrm{d}}

\begin{document}
\title{cs224n a3 mzx}
\maketitle

\section*{(1)}
\subsection*{a}
\noindent \textbf{\romannumeral1}
one batch only count $1-\beta_1$ (usually $10\%$) in the update, so the update is less likely to swing the direction. so we can get a more stable convergence

\noindent \textbf{\romannumeral2}
$m/\sqrt{v}$ will be normalised to a suitable size. parameters with small $v$ get larger updates (compared with pure sgd); params with small gradients can still have a suitable update speed. params with large gradients are scaled down can reduce the oscillation.

\subsection*{b}

\noindent \textbf{\romannumeral1}
\[
\gamma = \frac{1}{1- p_{\mathrm{drop}}}
\]

\begin{align*}
\mathbb{E}_{p_{\mathrm{drop}}}[h_{\mathrm{drop}}]_i
&= \mathbb{E}_{p_{\mathrm{drop}}}[\gamma d \odot h]_i\\
&= \mathbb{E}_{p_{\mathrm{drop}}} [\frac{1}{1- p_{\mathrm{drop}}} d \odot h]_i\\
&= \frac{1}{1- p_{\mathrm{drop}}}  \mathbb{E}_{p_{\mathrm{drop}}} [d \odot h]_i \\
&= \frac{1}{1- p_{\mathrm{drop}}} h_i  \mathbb{E}_{p_{\mathrm{drop}}}[d_i]\\
&= \frac{1}{1- p_{\mathrm{drop}}}  h_i  (1- p_{\mathrm{drop}}) \\
&= h_i
\end{align*}
\noindent \textbf{\romannumeral2}

when training dropout is used to prevent overfitting, help all the neurons learning usfull info\\

at evaluation we want the model's full capacity and a deterministic output, so we use all units with no random dropping; otherwise the same input could give different
predictions. the $\gamma$ scaling already matches the expected activation, so at test time we just use the full $h$ directly.

\section*{(2)}
\subsection*{a}

\begin{tabular}{llll}
\hline
stack & buffer & new dependency & transition \\
\hline
{[ROOT]} & [I, parsed, this, sentence, correctly] & & Initial Configuration \\
{[ROOT, I]} & [parsed, this, sentence, correctly] & & SHIFT \\
{[ROOT, I, parsed]} & [this, sentence, correctly] & & SHIFT \\
{[ROOT, parsed]} & [this, sentence, correctly] & parsed$\to$I & LEFT-ARC \\
% mine
{[ROOT, parsed, this]} & [sentence, correctly] & & SHIFT \\
{[ROOT, parsed, this, sentence]} & [ correctly] & & SHIFT \\
{[ROOT, parsed,  sentence]} & [ correctly] & sentence$\to$this&  LEFT-ARC \\
{[ROOT, parsed]} & [ correctly] & parsed$\to$sentence&  RIGHT-ARC \\
{[ROOT, parsed, correctly]} & [ ] & & SHIFT \\
{[ROOT, parsed]} & [ ] & parsed$\to$correctly & RIGHT-ARC \\
{[ROOT]} & [ ] & ROOT$\to$parsed & RIGHT-ARC \\
\hline
\end{tabular}
\subsection*{b}
$2n$ each of the n words is shifted from buffer to stack once, so n shifts. And each of the n words is removed by an arc action, so n times of arc. in total $2n$
\subsection*{e}
Best dev UAS: 88.75; 
Test UAS: 88.97
\subsection*{f}
\noindent \textbf{\romannumeral1}\\
Error type: Verb Phrase Attachment Error \\
Incorrect dependency: wedding $\to$ fearing \\
Correct dependency: heading $\to$ fearing \\

\noindent \textbf{\romannumeral2}\\
Error type: Coordination Attachment Error \\
Incorrect dependency: makes $\to$ rescue \\
Correct dependency: rush $\to$ rescue \\
\noindent \textbf{\romannumeral3}\\
Error type: Prepositional Phrase Attachment Error \\
Incorrect dependency: named $\to$ Midland \\
Correct dependency: Joe $\to$ Midland \\

\noindent \textbf{\romannumeral4}\\
Error type: Modifier Attachment Error \\
Incorrect dependency: elements $\to$ most \\
Correct dependency: crucial $\to$ most \\

\end{document}